\chapter{Пространственная \texorpdfstring{$SO(3)$}{SO(3)}-суперсвязность
и предел следа \texorpdfstring{$M_{35}$}{M35}}

Предыдущий гейт установил совместимость двух требований: гладкий
монопольный профиль может иметь конечную энергию и индекс Каллиаса один,
причём оба результата управляются тем же граничным хопфовым классом.
Теперь требуется более строгая проверка: следует ли этот профиль из
связывающего родителя $M_{35}$ или был незаметно импортирован извне.

Проверка проводится до феноменологии. Ни константа связи, ни масштаб
вакуумного поля, ни новый калибровочный сектор не вводятся.

\section{Что действительно фиксирует след
\texorpdfstring{$M_{35}$}{M35}}

Связывающая алгебра
\begin{equation}
 \mathcal L(E)=
 \begin{pmatrix}M_{20}&E\\E^*&M_{15}\end{pmatrix}
 \simeq M_{35}(\mathbb C)
\end{equation}
имеет один нормированный след. Он канонически фиксирует веса двух углов
$4/7$ и $3/7$, норму относительной внутренней кривизны и сокращение
хопфовых линий в диагональных произведениях.

Но конечный след действует только на матричных индексах. Для скалярной
координатной функции $f(x)$ выполнено
\begin{equation}
 [f(x)I_{35},X]=0,
 \qquad X\in M_{35}(\mathbb C).
 \label{eq:v5-spatial-trace-no-derivative}
\end{equation}
Следовательно, один $M_{35}$ не порождает $\partial_i f$, пространственный
оператор Дирака, внешнее произведение форм или звезду Ходжа. Для поля
$A_i(x)$ необходимо дополнительно задать по меньшей мере произведение с
пространственным дифференциальным исчислением.

Это не запрещает использовать уже существующее физическое пространство.
Но различает два утверждения:
\begin{equation}
 \text{совместимость с пространственным исчислением}
 \neq
 \text{вывод исчисления из }M_{35}.
\end{equation}

\section{След не выбирает калибровочную подалгебру}

Если флуктуировать весь связывающий контейнер, его унитарная алгебра имеет
размерность
\begin{equation}
 \dim\mathfrak u(35)=1225.
\end{equation}
Даже сохранение двух углов оставляет
\begin{equation}
 \dim\bigl(\mathfrak u(20)\oplus\mathfrak u(15)\bigr)=625.
\end{equation}
Нужная семейная алгебра $\mathfrak{so}(3)$ имеет размерность три.
Нормированный след инвариантен относительно всех внутренних унитарных
преобразований и потому сам не отличает эту трёхмерную подалгебру от
остальных направлений.

Прежний запрет остаётся в силе: объявить весь $M_{35}$ координатной
алгеброй и калибровать $U(35)$ нельзя. Допустимая $SO(3)$-связность должна
быть выбрана уже существующим семейным представлением или новым
родительским ограничением, но не одним следом.

\begin{nogocriterion}[След не выбирает группу]
Положительная инвариантная форма $\tau_{35}(X^*X)$ нормирует уже выбранную
калибровочную алгебру, но не выбирает саму подалгебру
$\mathfrak{so}(3)\subset\mathfrak u(35)$.
\end{nogocriterion}

\section{Топологическое препятствие гладкого ядра}

На сфере вокруг дефекта ориентированная стрелка несёт линию $L$ с
\begin{equation}
 c_1(L)=1.
\end{equation}
Однако
\begin{equation}
 H^2(B^3,\mathbb Z)=0.
\end{equation}
Ограничение любой комплексной линии на шаре $B^3$ к его границе имеет
нулевой первый класс Черна. Поэтому $L$ не может продолжиться через ядро
как всюду определённая невырожденная линия.

В гладком монополе это препятствие снимается иначе. Линия $L$ является
положительным собственным подрасслоением ранга-два пакета. В центре поле
$\Phi$ обращается в нуль, разложение на $L$ и $L^*$ перестаёт быть
определённым, а полный пакет
\begin{equation}
 L\oplus L^*,
 \qquad c_1(L\oplus L^*)=0,
 \label{eq:v5-hopf-pair-trivial-total-chern}
\end{equation}
гладко продолжается через шар.

Следовательно, конечная энергия требует не только граничной линии, но и
комплексно-линейного смешивания двух ветвей на ядре. Антилинейная
вещественная структура KO6, обменивающая $L$ и $L^*$, сама по себе не
является таким калибровочным смешиванием.

\section{Возврат представительного препятствия}

На частичной семейной половине текущий модуль раскладывается как
\begin{equation}
 \mathbf1\oplus\mathbf3\oplus\mathbf3\oplus\mathbf3,
\end{equation}
то есть содержит спины $0,1,1,1$. Для собственного расслоения состояния
с максимальным весом $j$ над $S^2$ класс Черна равен
\begin{equation}
 c_1=2j.
 \label{eq:v5-spin-j-chern}
\end{equation}
Поэтому текущие нетривиальные триплеты дают класс два, а не один:
\begin{equation}
 (0,2,2,2).
\end{equation}

Единичная линия требует полуцелого представления $j=1/2$. Но предыдущие
гейты уже доказали, что такого комплексного дублета нет в $H_{15}/M_{35}$:
KO6-пара обменивается антилинейно, слабый дублет имеет иные квантовые
числа, а новый семейный дублет меняет бюджет состояний и след.

Хопфов ориентационный функтор не отменяет этот запрет. Он канонически
назначает линию стрелке на границе, но не создаёт комплексно-линейные
операторы, смешивающие $L$ и $L^*$ в гладком ядре.

\begin{theorem}[Граничная ориентация не равна гладкому дублету]
Родитель $M_{35}$ канонически фиксирует хопфову линию с $c_1=1$ на сфере
вокруг дефекта, но текущий семейный модуль содержит только целоспиновые
представления. Поэтому он не реализует гладкий ранга-два $SU(2)$-пакет,
внутри которого эта линия могла бы исчезнуть как отдельное собственное
подрасслоение на ядре.
\end{theorem}

\section{Проверка квадрата суперсвязности}

Если внешнее пространственное исчисление, $SO(3)$-вложение и поле $\Phi$
уже заданы, суперсвязность
\begin{equation}
 \mathbb A=\nabla_A+\Phi
\end{equation}
имеет кривизну
\begin{equation}
 \mathbb F=F_A+D_A\Phi+\Phi^2.
\end{equation}
Один след действительно упаковывает калибровочную, кинетическую и
потенциальную части. В этом ограниченном смысле следовая архитектура
совместима с монопольным действием.

Но совместимость ещё не фиксирует отображение физического поля во
внедиагональный угол. При замене
\begin{equation}
 \Phi\longmapsto a\Phi
\end{equation}
схематические коэффициенты становятся
\begin{equation}
 |F_A|^2+2a^2|D_A\Phi|^2+a^4|\Phi|^4+\cdots.
 \label{eq:v5-superconnection-rescaling}
\end{equation}
После объявления вложения коэффициенты фиксируются одним следом, но сам
масштаб $a$, вакуум $v$ и пространственная звезда Ходжа из $M_{35}$ не
следуют. Предыдущий гейт суперсвязности уже обнаружил аналогичную границу:
общий множитель следа фиксируется, а относительное кинетико-потенциальное
действие требует дополнительной эрмитовой структуры.

Возможность получить действие Янга--Миллса--Хиггса из суперсвязности
известна математически; она подтверждает допустимость языка, но не выбирает
наши проектные данные. См.
\path{arXiv:hep-th/9701179} и
\path{arXiv:hep-th/0310267}.

\section{Сопоставление с ранними результатами}

Гейт \path{family_connection_defect_gap_bridge.tex} остаётся верным: на
трубчатой окрестности можно задать вещественную $SO(3)$-связность и
получить точное одномерное ядро. Но эта связность была условным входом и
не имела плоского продолжения по торсионному циклу с требуемым углом.

Нынешняя хопфова линия улучшает граничную топологию и ориентацию, но не
выводит глобальную пространственную связность. Таким образом, два
результата согласуются:
\begin{equation}
 \text{локальная }SO(3)\text{-кинематика существует},
 \qquad
 \text{родительская динамика не выведена}.
\end{equation}

\section{Вердикт}

Тест следа $M_{35}$ даёт смешанный, но однозначный ответ:
\begin{equation}
 \boxed{
 \begin{gathered}
 \text{граничная хопфова ориентация: выведена},\\
 \text{суперсвязностная упаковка после задания полей: допустима},\\
 \text{пространственное исчисление и }SO(3)\text{-вложение: не выведены},\\
 \text{гладкий единично-индексный дублет в }M_{35}:\text{ отсутствует}.
 \end{gathered}}
\end{equation}

Поэтому стандартное BPS-завершение остаётся математическим образцом, но
не является действием версии V.

Остаётся один действительно новый внутренний вопрос. Пара
$E\otimes L$ и $E^*\otimes L^*$ уже присутствует как две степени
градуированного соответствия. Возможно, гладкое смешивание $L/L^*$ должно
нести не чётная калибровочная связность, а нечётная компонента самой
суперсвязности, обращающаяся в нуль на дефекте. Это не равносильно новому
фиксированному $SU(2)_F$-дублету и потому заслуживает отдельного теста.

Следующий гейт:
\path{version5_hopf_pair_odd_core_extension_gate.tex}.

\begin{protocol}
\begin{enumerate}
 \item один внутренний след $M_{35}$: \textbf{есть};
 \item пространственная производная из конечного следа: \textbf{нет};
 \item выбор $\mathfrak{so}(3)\subset\mathfrak u(35)$ следом:
 \textbf{нет};
 \item чистая линия $L$ продолжается через $B^3$: \textbf{нет};
 \item пара $L\oplus L^*$ топологически продолжается: \textbf{да};
 \item комплексно-линейное смешивание пары в текущем модуле:
 \textbf{не получено};
 \item классы Черна текущих семейных представлений:
 \textbf{нулевые или чётные};
 \item единичный класс из текущего семейного триплета: \textbf{нет};
 \item один след после внешнего задания суперсвязности:
 \textbf{совместим};
 \item BPS-веса и масштаб из одного $M_{35}$: \textbf{не выведены};
 \item физическое замыкание: \textbf{не достигнуто}.
\end{enumerate}
\end{protocol}

Машинный сертификат сохранён в
\path{s2t_v5_spatial_so3_superconnection_parent_trace_gate_results.json}.